% Generated draft: requires mathematical and source review.
\documentclass[11pt]{article}
\usepackage[margin=25mm]{geometry}
\usepackage{fontspec}
\setmainfont{Latin Modern Roman}
\usepackage{amsmath,amssymb,mathtools,unicode-math}
\setmathfont{Latin Modern Math}
\usepackage{enumitem,needspace,xurl,hyperref}
\setlength{\parindent}{0pt}
\setlength{\parskip}{5pt}
\setlength{\emergencystretch}{3em}
\newcommand{\N}{\mathbb{N}}
\newcommand{\Z}{\mathbb{Z}}
\newcommand{\Q}{\mathbb{Q}}
\newcommand{\R}{\mathbb{R}}
\newcommand{\C}{\mathbb{C}}
\newcommand{\F}{\mathbb{F}}
\newcommand{\A}{\mathbb{A}}
\newcommand{\PP}{\mathbb{P}}
\newcommand{\E}{\mathbb{E}}
\newcommand{\eps}{\varepsilon}
\newcommand{\dd}{\,\mathrm{d}}
\begin{document}
\section*{UnsolvedMath ID: 800553}
\textbf{Source problem code:} catalogue entry
\subsection{Definitions and mathematical statement}
For a finite graph \(G\), let \(\mathcal D(G)\) be the set of plane drawings in which vertices are distinct points and edges are simple arcs joining their endpoints, with edge interiors avoiding vertices. Use drawings in general position: interior intersections are transverse crossings, there are finitely many crossings, and no three edge interiors meet at one point. The crossing number is
\[
\operatorname{cr}(G)=\min_{D\in\mathcal D(G)}\#\{\text{interior edge crossings in }D\}.
\]
This uses the usual finite-graph drawing convention; the supplied record does not explicitly specify the graph category.

Let \(\mathcal B_k\) denote the class called \emph{bar \(k\)-visibility graphs} in the source. The requested extremal quantity is
\[
\max\{\operatorname{cr}(G):G\in\mathcal B_k\},
\]
provided that this maximum exists. The question is to determine the largest crossing number in this class.

The supplied record does not define bar \(k\)-visibility, specify the admissible values of \(k\), or impose a number of vertices or another size bound. Thus \(\mathcal B_k\) remains a source-dependent, unresolved part of the formulation. In particular, no fixed-order restriction may be silently inserted, and existence or finiteness of the displayed maximum is not asserted. A fully specified extremal problem requires clarification of these points.
\subsection{Short English statement}
Determine the largest possible crossing number of a bar \(k\)-visibility graph. The supplied record leaves the visibility model, the range of \(k\), and any intended graph-size restriction unspecified, so the exact extremal target requires clarification.
\subsection{Sources}
\begin{itemize}
\item Local UnsolvedMath registry, supplied complete source record: permanent ID \(800553\), exact source code \texttt{catalogue entry}, title \emph{Visibility Graphs of Points, Segments and Polygons --- Problem 27}. The registry attributes the problem to Ghosh and Goswami, with proposed year 2010, and labels it open; this status is supplied metadata, not an independently verified mathematical assessment.
\item Ghosh and Goswami, \emph{Unsolved Problems in Visibility Graphs of Points, Segments and Polygons} (2010), Problem 27. Supplied source URL: \url{https://arxiv.org/abs/1012.5187}. The record reports access on 2026-10-16 and extraction by \texttt{arxiv\_latex}; these are supplied provenance details, not a claim of independent consultation or verification.
\end{itemize}
\end{document}
